\subsection{Continual Training Setup}
\label{app:continue_training}
In Section~\ref{sec:robustness}, we introduce the continual training method \grit on StackExchange data to evaluate whether training on in-domain data enhances the performance of \ours. Detailed experimental settings are described in this section. Specifically, we follow \grit to train models with two distinct objectives: a contrastive loss to maintain the model's retrieval capability and a language modeling loss to preserve the model's language generation ability. For training with the contrastive loss, we collect 3,200 (post, answer) pairs from the \biology, \earth, \econ, \psychology, \robotics, and \stackoverflow sections of StackExchange, and 1,538 pairs from \sustainable. Each post's answer is used as a positive example, with other answers serving as in-batch negatives. 
For training with the language modeling loss, we use both positive and negative documents from each domain within the StackExchange subsection of \ours{}. These documents are split into chunks of 2048 tokens, and we sample up to 3,200 chunks for training. We use a small batch size of 64 to ensure sufficient learning steps, while following the other hyperparameters as outlined in ~\citet{muennighoff2024generative}. We continue training \grit for 10 epochs, benchmarking the checkpoint from each epoch on the StackExchange datasets of \ours. The detailed scores are in Table~\ref{tab:robustness_details}, where we copy the scores of \grit to epoch=0 for easier reference. The results indicate no significant improvement across the 10 epochs, suggesting that even with intensive inclusion of StackExchange data or relevant domain knowledge in the training data of language models or retrievers, performance may not increase substantially without enhancing incorporating reasoning into the retrieval process.


\begin{table}[htbp]
\centering
\caption{\textbf{Scores of finetuned \grit of every epoch on StackExchange datasets of \ours.} Epoch=0 indicates the performance of \grit without further training.}
% \resizebox{\textwidth}{!}{
\begin{tabular}{l|ccccccc|c}
\toprule
\textbf{Epoch} & \textbf{Bio.} & \textbf{Earth.} & \textbf{Econ.} & \textbf{Psy.} & \textbf{Rob.} & \textbf{Stack.} & \textbf{Sus.} & \textbf{Avg.} \\
\midrule
0 (GritLM) & 25.0 & 32.8 & 19.0 & 19.9 & 17.3 & 11.6 & 18.0 & 20.5\\
1 & 22.2 & 25.4 & 17.6 & 28.1 & 11.1 & 9.8 & 19.6 & 19.1 \\
2 & 18.7 & 23.8 & 13.5 & 19.3 & 10.7 & 10.2 & 16.5 & 16.1 \\
3 & 20.9 & 23.6 & 16.9 & 25.2 & 11.1 & 8.5 & 16.6 & 17.5 \\
4 & 24.3 & 28.0 & 18.3 & 26.9 & 13.4 & 13.3 & 20.0 & 20.6 \\
5 & 23.1 & 28.5 & 18.4 & 26.1 & 14.6 & 11.7 & 21.6 & 20.6 \\
6 & 19.9 & 26.4 & 16.0 & 27.9 & 9.6 & 9.3 & 19.3 & 18.3 \\
7 & 24.3 & 25.4 & 16.5 & 28.1 & 11.0 & 9.8 & 17.0 & 18.9 \\
8 & 21.6 & 28.7 & 19.2 & 28.7 & 11.1 & 11.8 & 22.4 & 20.5 \\
9 & 21.3 & 29.0 & 20.0 & 28.7 & 11.4 & 14.3 & 22.0 & 21.0 \\
10 & 21.1 & 25.5 & 18.8 & 30.7 & 12.7 & 12.1 & 21.9 & 20.4\\
\bottomrule
\end{tabular}
% }
\label{tab:robustness_details}
\end{table}


